% !TEX root = document.tex

\chapter{\label{chap:introduction}Introduction}

Language workbenches aim to reduce the cost of designing and implementing programming languages by providing high-level, declarative abstractions for syntax, name resolution, and type checking. Spoofax \autocite{6898704}, developed at TU Delft, is one such workbench. It enables users to define the syntax of a language through a SDF3 grammar specification \autocite{sdf3}, from which an abstract syntax tree (AST) representation is automatically derived. This AST forms the basis for subsequent analyses, including name resolution and type checking.

Type checking in Spoofax is specified using Statix \autocite{van_antwerpen_scopes_2018}, a constraint-based language designed for expressing typing and name resolution rules declaratively. Statix combines term constraints with scope graphs, a graph-based representation of scopes and the relations between them. Scope graphs are particularily well suited for imports, mutual refferences and other types of non-lexical scoping.

At a conceptual level, Statix bears similarities to $\lambda$-Prolog, a higher-order logic programming language that extends Prolog with types and higher-order terms. In $\lambda$-Prolog, type systems can be encoded directly as logical specifications using predicates, quantification, and higher-order unification. Terms in $\lambda$-Prolog are written using higher order abstract syntax (HOAS) which makes the language particularly well-suited for turning formal type checker rules into an implementation.

HOAS and Scope Graphs differ fundamentally in how they treat binding and unification. $\lambda$-Prolog relies on meta-level quantification and higher-order unification to model binding, whereas Statix represents binding explicitly using scope graphs and resolves names via constraints. While scope graphs provide strong support for modular and non-lexical scoping, they do not natively provide higher-order pattern unification, which is central to many dependently typed systems.

This tension motivates the central question of this thesis: can higher-order pattern unification be integrated with a scope-graph-based constraint system in a principled and practical way? Furthermore, how does such an approach compare to higher-order logic programming systems when implementing a dependently typed language?

To address these questions, we design and evaluate a higher-order pattern unification algorithm tailored to a scope-graph-based setting, and analyze its implications for implementing dependently typed languages in Spoofax.

\section{\label{sec:contributions}Contributions}

This thesis addresses two research questions and makes the following contributions.

\paragraph{Higher-order pattern unification in scope-graph systems (RQ1).}
We design and formalize an algorithm for higher-order pattern unification in a constraint-based system that uses scope graphs for name binding. To validate the algorithm, we implement a minimal dependently typed core language in Statix, demonstrating that higher-order pattern unification can be integrated with scope-graph-based name resolution and supports core dependent type features.

\paragraph{Comparison of implementation paradigms for dependent typing (RQ2).}
We provide a structured comparison between higher-order logic programming systems such as $\lambda$-Prolog and scope-graph-based constraint systems. The analysis examines differences in expressivity and computability, including the treatment of binding, fresh name generation, modularity, and non-lexical scoping. This clarifies the trade-offs between meta-level binding mechanisms and explicit scope representations when implementing practical dependently typed languages.